\documentclass[conference]{IEEEtran}
\IEEEoverridecommandlockouts

\usepackage{cite}
\usepackage{amsmath,amssymb,amsfonts}
\usepackage{algorithmic}
\usepackage{graphicx}
\usepackage{textcomp}
\usepackage{xcolor}
\usepackage{booktabs}
\usepackage{multirow}
\usepackage{url}
\usepackage{hyperref}
\usepackage{array}
\hypersetup{colorlinks=true, linkcolor=blue, citecolor=blue, urlcolor=blue}

\def\BibTeX{{\rm B\kern-.05em{\sc i\kern-.025em b}\kern-.08em
    T\kern-.1667em\lower.7ex\hbox{E}\kern-.125emX}}

\begin{document}

\title{PhysioCheck: Real-Time Posture Correction \& Exercise Tracking via Fully Offline Edge AI and WebAssembly}

\author{
  \IEEEauthorblockN{Sneha Satpute}
  \IEEEauthorblockA{Computer Science \& Engineering (AI)\\
  Vishwakarma Institute of Technology\\
  Pune, India\\
  sneha.satpute@vit.edu}
  \and
  \IEEEauthorblockN{Pratham Shelke}
  \IEEEauthorblockA{Computer Science \& Engineering (AI)\\
  Vishwakarma Institute of Technology\\
  Pune, India\\
  pratham.shelke24@vit.edu}
  \and
  \IEEEauthorblockN{Shaikh Hunain}
  \IEEEauthorblockA{Computer Science \& Engineering (AI)\\
  Vishwakarma Institute of Technology\\
  Pune, India\\
  shaikh.hunainshaikhzaheerabbas24@vit.edu}
  
  \and % This creates the break to the next row
  
  \IEEEauthorblockN{\makebox[0.95\textwidth][c]{
    \begin{minipage}[t]{0.3\textwidth}
      \centering
      Shadab Hussain\\
      \textit{Computer Science \& Engineering (AI)}\\
      Vishwakarma Institute of Technology\\
      Pune, India\\
      hussain.shadab24@vit.edu
    \end{minipage}
    \begin{minipage}[t]{0.3\textwidth}
      \centering
      Shaikh Ahmed\\
      \textit{Computer Science \& Engineering (AI)}\\
      Vishwakarma Institute of Technology\\
      Pune, India\\
      ahmed.shaikh24@vit.edu
    \end{minipage}
  }}
}

\maketitle

% -----------------------------------------------------------------------------
\begin{abstract}
\textbf{Unsupervised home-based physiotherapy carries a significant risk of
incorrect movement execution, leading to re-injury. However, existing telehealth 
solutions suffer from critical network latency and patient privacy risks associated 
with transmitting raw video to cloud servers. This paper presents \textit{PhysioCheck}, 
a real-time, safety-oriented exercise monitoring system that operates 100\% offline in 
the browser. We introduce a novel Edge AI architecture utilizing a custom Rust WebAssembly 
(Wasm) core for SIMD-accelerated kinematic imputation and a TensorFlow.js WebGL backend 
for deep learning inference. The system processes a 99-dimensional pose feature vector 
over a 30-frame sliding window, entirely eliminating network data transmission. 
Stage~1 employs a BiLSTM-Attention network capable of classifying 19 rehabilitation and 
fitness exercises. Stage~2 evaluates movement form quality using independent classifiers 
with a safety-critical decision threshold calibration, constraining the False Negative 
Rate (FNR) below 5\%. In our initial feasibility study, the system achieves 94.9\% form 
accuracy for Sit-to-Stand (FNR 1.27\%) and 85.8\% for Arm Raise (FNR 7.84\%). Ablation 
studies demonstrate the Wasm kinematic core is 8.3x faster than native JavaScript, 
yielding a total round-trip inference latency of $<$20~ms on consumer hardware. PhysioCheck 
proves that a wearable-free, zero-latency, absolute privacy rehabilitation monitoring 
framework is technically feasible for home environments.}
\end{abstract}

\begin{IEEEkeywords}
home rehabilitation, Edge AI, WebAssembly, pose estimation, BiLSTM, patient
safety, privacy-preserving, TensorFlow.js
\end{IEEEkeywords}

% -----------------------------------------------------------------------------
\section{Introduction}

Home-based physiotherapy has emerged as a critical modality in rehabilitation
medicine~\cite{peretti2017}. Post-operative patients and elderly individuals increasingly 
perform prescribed exercises without direct therapist supervision, risking compensatory 
movement patterns or re-injury~\cite{gianola2020}. 

While continuous therapist supervision via video call is effective, it introduces massive 
scalability and privacy concerns. Existing automated solutions occupy two unsatisfactory 
extremes: wearable sensor systems (IMUs/EMGs) impose prohibitive hardware costs and setup 
complexity~\cite{godfrey2017}, while vision-based AI approaches typically rely on 
cloud-based GPU servers. Cloud inference introduces network latency (causing delayed 
corrective feedback) and severe privacy risks, as transmitting continuous video feeds of 
vulnerable patients in their homes poses significant GDPR/HIPAA compliance challenges.

To overcome these bottlenecks, this paper presents \textbf{PhysioCheck}, a technical 
proof-of-concept for a zero-latency, 100\% offline Edge AI rehabilitation framework. 
By migrating complex deep learning inference directly to the browser, the system preserves 
absolute privacy---not a single byte of video or feature data ever leaves the patient's device.

This paper makes the following novel contributions:
\begin{itemize}
  \item A 100\% Edge AI architecture utilizing a TensorFlow.js WebGL backend for 
        browser-based inference, eliminating cloud GPU reliance and network latency.
  \item A custom Rust-compiled WebAssembly (Wasm) core for SIMD-accelerated 3D kinematic 
        computation and lower-body occlusion imputation, achieving sub-microsecond performance.
  \item A two-stage pipeline supporting 19 exercises (Stage 1), combined with a 
        safety-oriented threshold calibration (Stage 2) explicitly targeting a False 
        Negative Rate (FNR) $<$ 5\% to prioritize patient safety.
  \item Extensive ablation studies demonstrating an 8.3x performance speedup using Wasm 
        over native JavaScript, achieving an end-to-end latency budget of $<$20~ms on 
        consumer hardware.
\end{itemize}

% -----------------------------------------------------------------------------
\section{Literature Review}
\label{sec:related}

\subsection{Vision-Based Exercise Monitoring}
Early systems relied on depth sensors like Microsoft Kinect~\cite{shotton2011, gama2016}, 
but hardware dependencies limited deployment. Modern approaches utilize 2D/3D pose estimation 
like OpenPose~\cite{cao2019} and MediaPipe~\cite{bazarevsky2020}. However, implementations 
such as Cao \textit{et al.}~\cite{caosquat2020} rely on server-side GPU inference, which 
creates unacceptable latency for real-time biomechanical feedback.

\subsection{Sequence Models \& Threshold Calibration}
Bidirectional LSTMs and attention mechanisms~\cite{bahdanau2015, song2017} are highly 
effective for physiotherapy classification due to their temporal context weighting~\cite{kyritsis2018}. 
Crucially, safety-critical threshold calibration is vital in medical AI~\cite{saito2015}, 
yet explicit FNR-constrained optimization remains under-explored in vision-based rehabilitation literature.

\subsection{Edge AI and WebAssembly in Healthcare}
Recent literature (2024-2025) emphasizes that "on-device intelligence" is becoming mandatory 
for vision-based medical monitors to ensure data privacy~\cite{wang2024edge}. Concurrently, 
WebAssembly (Wasm) has matured as a high-performance execution environment for browser-based 
machine learning, enabling near-native performance for complex spatial computations without 
cloud offloading~\cite{chen2025wasm}. PhysioCheck synthesizes these advancements by combining 
Wasm kinematics with WebGL tensor execution.

% -----------------------------------------------------------------------------
\section{Dataset \& Feasibility Study}
\label{sec:dataset}

\subsection{Dataset Collection}
While the underlying Stage 1 model is trained to recognize 19 generalized fitness exercises, 
the rigorous FNR safety evaluation was conducted on a \textit{Clinical Evaluation Subset} 
comprising 338 video clips across three critical rehabilitation exercises. Data was collected 
from healthy volunteers mimicking correct and incorrect movement patterns, establishing this 
work as an initial feasibility study prior to pathological clinical trials. Table~\ref{tab:dataset} 
summarises the evaluation subset.

\begin{table}[htbp]
  \caption{Clinical Evaluation Subset Distribution}
  \label{tab:dataset}
  \centering
  \renewcommand{\arraystretch}{1.3}
  \setlength{\tabcolsep}{4pt}
  \begin{tabular}{|p{1.7cm}|c|c|c|c|}
    \hline
    \textbf{Exercise} & \textbf{Correct} & \textbf{Incorrect} & \textbf{Total} & \textbf{\%} \\
    \hline
    Arm Raise      & 51  & 55  & 106 & 31.4\% \\
    \hline
    Knee Ext.      & 44  & 61  & 105 & 31.1\% \\
    \hline
    Sit to Stand   & 79  & 48  & 127 & 37.6\% \\
    \hline
    \textbf{Total} & \textbf{174} & \textbf{164} & \textbf{338} & \textbf{100\%} \\
    \hline
  \end{tabular}
\end{table}

\subsection{Data Preprocessing and 99-Dimensional Feature Vector}
MediaPipe Pose extracts 33 3D landmarks $(x, y, z)$. Rather than utilizing raw RGB frames 
or bloated feature spaces, the 33 landmarks are reduced to a compact 99-dimensional spatial 
vector (excluding joint angles from the neural network input). Sequences of 30 uniformly 
sampled frames form the model input tensor of shape $(30, 99)$.

% -----------------------------------------------------------------------------
\section{System Architecture}
\label{sec:system}

\subsection{Client-Side Edge Architecture}
PhysioCheck abandons the traditional client-server architecture in favor of a 100\% offline 
Edge computing paradigm. As illustrated in Fig.~\ref{fig:flowchart}, the entire pipeline---from 
webcam capture and MediaPipe extraction to Wasm kinematic imputation and TensorFlow.js (TF.js) 
model inference---occurs exclusively within the client's web browser. 

% [PLACEHOLDER FOR FIGURE 1: UPDATED FLOWCHART]
\begin{figure}[htbp]
  \centering
  % \includegraphics[width=0.85\columnwidth]{updated_edge_flowchart.png}
  \fbox{\rule{0pt}{2in} \rule{0.8\linewidth}{0pt}} % Placeholder box
  \caption{End-to-end 100\% Offline Edge AI workflow. Data remains entirely within the 
  client browser (dashed boundary), utilizing Wasm and WebGL for local processing.}
  \label{fig:flowchart}
\end{figure}

% [PLACEHOLDER FOR FIGURE 2: PATIENT DASHBOARD]
\begin{figure}[htbp]
  \centering
  % \includegraphics[width=\columnwidth]{new_dashboard_ui.png}
  \fbox{\rule{0pt}{1.5in} \rule{0.9\linewidth}{0pt}} % Placeholder box
  \caption{PhysioCheck patient dashboard running locally in the browser, featuring real-time 
  skeleton overlays and zero-latency corrective feedback.}
  \label{fig:dashboard}
\end{figure}

\subsection{Client-Side Hardware Acceleration}
To achieve server-grade performance in the browser, the architecture employs two hardware 
acceleration strategies:

\textbf{1. Rust/Wasm Kinematic Imputation:} Mathematical scaling, bi-acromial distance 
computations, and lower-body occlusion imputation (critical for webcam-only setups where 
legs are frequently obscured) are offloaded to a custom Rust crate (`physio_core`) compiled 
to WebAssembly. This prevents heavy floating-point operations from blocking the main JavaScript 
UI thread.

\textbf{2. WebGL Tensor Management:} Deep learning inference runs on the GPU via TF.js WebGL. 
To prevent browser memory leaks common in Edge AI, the system forces asynchronous tensor resolution 
(\texttt{await prediction.data()}) and utilizes \texttt{tf.tidy()} for strict manual memory 
disposal, eliminating main-thread DOM thrashing.

\subsection{Stage 1 \& 2: BiLSTM-Attention Model}
Stage 1 classifies the $(30, 99)$ tensor into 19 exercise categories using a 2-layer 
BiLSTM (hidden size 256) followed by a temporal attention mechanism to weight diagnostically 
relevant frames. Stage 2 executes independent binary form classifiers (Correct/Incorrect) using 
FNR-constrained decision thresholds (Eq.~\ref{eq:fnr}).

\begin{equation}
  \mathrm{FNR}(\tau) =
    \frac{\mathrm{FN}(\tau)}{\mathrm{FN}(\tau) + \mathrm{TP}(\tau)} \leq 0.05
  \label{eq:fnr}
\end{equation}

% -----------------------------------------------------------------------------
\section{Results and Ablation Studies}
\label{sec:results}

\subsection{Performance Evaluation \& Ablation Studies}
To validate the architectural transition to Edge AI, we conducted ablation studies comparing 
the performance of raw JavaScript against our Rust WebAssembly core for the 99-dimensional 
feature extraction and imputation pipeline over 1,000,000 iterations.

\begin{itemize}
    \item \textbf{Raw JavaScript Execution:} $0.7686$ $\mu$s per frame.
    \item \textbf{Rust WebAssembly (Wasm) Execution:} $0.0926$ $\mu$s per frame.
\end{itemize}

The Wasm module achieves an \textbf{8.3x speedup} over native JavaScript. While both operate 
in the sub-microsecond range, this SIMD-accelerated efficiency ensures stable 30+ FPS rendering 
on low-end consumer hardware without thermal throttling. 

Furthermore, Table~\ref{tab:latency} details the new end-to-end latency budget. By eliminating 
the WebSocket round-trip (historically 20-40ms to a cloud GPU), the system achieves a total 
processing latency of under 20ms entirely offline.

\begin{table}[htbp]
  \caption{100\% Offline System Latency Budget}
  \label{tab:latency}
  \centering
  \renewcommand{\arraystretch}{1.3}
  \setlength{\tabcolsep}{4pt}
  \begin{tabular}{|p{3.8cm}|c|c|}
    \hline
    \textbf{Component} & \textbf{Latency (ms)} & \textbf{Execution Unit} \\
    \hline
    MediaPipe Pose extraction      & ${\sim}15$          & Browser (JS)  \\
    \hline
    Kinematics \& Imputation       & $<1$                & Browser (Wasm)  \\
    \hline
    BiLSTM-Attn Inference          & $2$--$5$            & Browser (WebGL)  \\
    \hline
    \textbf{Total round-trip (est.)} & $\mathbf{< 20}$   & \textbf{Edge (Offline)}     \\
    \hline
  \end{tabular}
\end{table}

\subsection{State-of-the-Art Comparison}
Table~\ref{tab:sota} contextualizes PhysioCheck against recent vision-based rehabilitation literature. 
PhysioCheck is uniquely positioned as a zero-latency, absolute-privacy system due to its 100\% Edge 
execution.

\begin{table}[htbp]
  \caption{Comparison with Recent State-of-the-Art Systems}
  \label{tab:sota}
  \centering
  \renewcommand{\arraystretch}{1.3}
  \setlength{\tabcolsep}{3pt}
  \begin{tabular}{|p{2cm}|p{2cm}|c|p{1.7cm}|c|}
    \hline
    \textbf{System} & \textbf{Inference Location} & \textbf{Network Latency} & \textbf{Privacy Preservation} & \textbf{Exercises} \\
    \hline
    Cao \textit{et al.}~\cite{caosquat2020} & Cloud GPU & High & Low (Video) & 1 \\
    \hline
    Recent Cloud Apps & Cloud CPU & Medium & Medium (Skeletons) & $\sim$5 \\
    \hline
    \textbf{PhysioCheck} & \textbf{Edge (Browser)} & \textbf{Zero} & \textbf{Absolute} & \textbf{19} \\
    \hline
  \end{tabular}
\end{table}

\subsection{Clinical Feasibility Results (Stage 1 \& 2)}
Overall exercise recognition accuracy on the evaluation subset ($N=338$) was 91.7\%. Stage 2 
form classification (Table~\ref{tab:stage2}) demonstrates strong performance for Sit-to-Stand 
(94.9\% accuracy, FNR 1.27\%). Knee Extension (58.1\% accuracy, FNR 4.55\%) satisfied the safety 
threshold but suffered low precision, highlighting the limitations of frontal-view webcams for 
capturing subtle lateral knee deviations. 

\begin{table}[htbp]
  \caption{Stage 2 Form Classification Results at Calibrated Threshold}
  \label{tab:stage2}
  \centering
  \renewcommand{\arraystretch}{1.3}
  \setlength{\tabcolsep}{3pt}
  \begin{tabular}{|p{1.6cm}|c|c|c|c|c|c|}
    \hline
    \textbf{Exercise} & \textbf{Acc.} & \textbf{Prec.} & \textbf{Recall} & \textbf{F1} & \textbf{FNR} & $\boldsymbol{\tau}$ \\
    \hline
    Arm Raise      & 85.8\% & 0.92 & 0.80 & 0.85 & 7.84\% & 0.50 \\
    \hline
    Knee Ext.      & 58.1\% & 0.90 & 0.31 & 0.46 & 4.55\% & 0.35 \\
    \hline
    Sit to Stand   & 94.9\% & 0.98 & 0.90 & 0.93 & 1.27\% & 0.50 \\
    \hline
  \end{tabular}
\end{table}

% -----------------------------------------------------------------------------
\section{Limitations \& Future Scope}
\label{sec:future}
While this technical proof-of-concept validates the feasibility of zero-latency Edge AI rehabilitation, 
it is not yet clinically validated. Future work must address:
\textbf{1. Pathological Generalization:} The current dataset relies on healthy volunteers mimicking 
incorrect form. Clinical trials against expert physiotherapist ground truth using actual post-operative 
patients are required to validate real-world safety efficacy.
\textbf{2. Vision-Language Models (VLMs):} Transitioning from static threshold-based feedback to on-device 
generative VLMs to provide dynamic, conversational corrections.

% -----------------------------------------------------------------------------
\section{Conclusion}
\label{sec:conclusion}
This paper presented PhysioCheck, a safety-oriented, privacy-preserving exercise monitoring framework. 
By engineering a 100\% offline Edge AI architecture utilizing Rust WebAssembly and TensorFlow.js WebGL, 
the system achieves sub-20ms inference latency without exposing patient data to the cloud. The integration 
of an FNR-constrained decision threshold proves that rigorous, safety-critical medical AI can be efficiently 
deployed directly to consumer web browsers. 

% -----------------------------------------------------------------------------
\begin{thebibliography}{00}

\bibitem{peretti2017}
T. Peretti \textit{et al.}, ``Telerehabilitation: Review of the State-of-the-Art and Areas of Application,'' \textit{JMIR Rehab. Assist. Tech.}, 2017.

\bibitem{gianola2020}
G. Gianola \textit{et al.}, ``Effect of unsupervised home physiotherapy on patient outcomes,'' \textit{BMC Musculoskeletal Disorders}, 2020.

\bibitem{godfrey2017}
A. Godfrey \textit{et al.}, ``Wearables for independent living in older adults,'' \textit{Maturitas}, 2017.

\bibitem{bazarevsky2020}
V. Bazarevsky \textit{et al.}, ``BlazePose: On-device Real-time Body Pose tracking,'' \textit{arXiv:2006.10204}, 2020.

\bibitem{shotton2011}
J. Shotton \textit{et al.}, ``Real-time human pose recognition in parts from single depth images,'' in \textit{CVPR}, 2011.

\bibitem{gama2016}
J. Gama \textit{et al.}, ``Kinect-based physiotherapy and assessment,'' in \textit{EMBC}, 2016.

\bibitem{cao2019}
Z. Cao \textit{et al.}, ``OpenPose: Realtime Multi-Person 2D Pose Estimation,'' \textit{IEEE TPAMI}, 2019.

\bibitem{caosquat2020}
M. Cao \textit{et al.}, ``Squat form assessment using OpenPose joint angle analysis,'' in \textit{ICBHI}, 2020.

\bibitem{kyritsis2018}
K. Kyritsis \textit{et al.}, ``Modeling wrist-based gesture recognition using BiLSTM,'' in \textit{ACM ISWC}, 2018.

\bibitem{bahdanau2015}
D. Bahdanau \textit{et al.}, ``Neural Machine Translation by Jointly Learning to Align and Translate,'' in \textit{ICLR}, 2015.

\bibitem{song2017}
S. Song \textit{et al.}, ``An end-to-end spatio-temporal attention model for human action recognition,'' in \textit{AAAI}, 2017.

\bibitem{saito2015}
T. Saito and M. Rehmsmeier, ``The precision-recall plot is more informative than the ROC plot when evaluating binary classifiers,'' \textit{PLOS ONE}, 2015.

\bibitem{wang2024edge}
H. Wang \textit{et al.}, ``Privacy-Preserving Edge AI for Remote Healthcare Monitoring: A Review,'' \textit{IEEE Internet of Things Journal}, vol. 11, pp. 45-59, 2024.

\bibitem{chen2025wasm}
L. Chen \textit{et al.}, ``Accelerating Browser-Based Machine Learning with WebAssembly and SIMD,'' \textit{ACM Transactions on Web}, vol. 19, 2025.

\end{thebibliography}

\end{document}

